\SourceLeaf{33} Mittheilung von Herrn \textsc{Mehmke} danke ich die Kenntniss, dass \textsc{Gauss} auch in der
zweiten Tabelle (abgesehen von den drei genannten Fällen) für das System
derjenigen Einheiten, welche mit Hülfe ganzzahliger $a$, $b$, $c$ in der
Gestalt $a+b\nu+c\nu\nu$ darstellbar sind, jedesmal die zur
Fundamentaleinheit reciproke Einheit angegeben hat. In der Mehrzahl der
Fälle handelt es sich auch um die Fundamentaleinheit des Körpers.
Ausnahmen von dieser Regel liegen vor bei $n=16$, wo die achte Potenz der
Fundamentaleinheit angegeben ist, und bei $n=10$ und $n=19$, wo die
reciproken Fundamentaleinheiten die folgenden sind:

\[
\frac{7+\sqrt[3]{10}-2\sqrt[3]{100}}{3}, \qquad
\frac{2+2\sqrt[3]{19}-\sqrt[3]{19^2}}{3}.
\]

\par\noindent\hfill\textsc{Fricke.}\par\endgroup
